\subsection{Fixed-dimensional likelihood localization}
\label{sec:support-likelihood-localization}

Fixed-dimensional likelihood localization is the first step in converting
the field estimates into the exact support identity of
Theorem~\ref{thm:main-support-law}. We apply the one-atom
geometry around a fit that already carries a small amount of mass away from
its central atom. The resulting comparison shows that edge mass cannot create
new KKT contacts where the central directional field has a quantitative gap
below one.

It is convenient to work in score coordinates. Replace a component location
\(u\) by \(t=\sqrt{1-\eps_n}\,u\); this bijection preserves support size.
We reuse \(G\) for the transformed mixing measure. Recall that
\(\gamma_n=(1-\eps_n)^{-1}=1+2\kappa_n\). If the central atom is at
\(\tau\), the likelihood ratio of the atom at \(\tau+a\) to the atom at
\(\tau\), evaluated at \(Z_i\), is
\[
    q_{n,\tau,a}(i)
    =\exp\left\{a\cdot(Z_i-\gamma_n\tau)
                    -\frac{\gamma_n|a|^2}{2}\right\}.
\]
For a perturbation \(G=(1-W)\delta_\tau+G^{\rm out}\), write
\[
    m_G(i)
    =1-W+\int q_{n,\tau,u-\tau}(i)\,G^{\rm out}(du)
\]
for its likelihood ratio to \(\delta_\tau\), and define
\[
    D_{n,\tau}(a)=\frac1n\sum_{i=1}^nq_{n,\tau,a}(i),
    \qquad
    \Psi_G(\tau+a)=\frac1n\sum_{i=1}^n
        \frac{q_{n,\tau,a}(i)}{m_G(i)}.
\]
The KKT conditions are \(\Psi_G\le1\), with equality at every support point.

The deterministic no-propagation estimate below separates the total mass of
the edge atoms from the geometry of their locations.
\begin{lemma}
\label{lem:support-no-percolation}
For every measure \(G=(1-W)\delta_\tau+G^{\rm out}\) defined above,
\[
    \Psi_G(\tau+a)\le \frac{D_{n,\tau}(a)}{1-W},
    \qquad a\in\mathbb R^d.
\]
Consequently, if \(D_{n,\tau}\le1-\eta\) on a set \(A\) and
\(W\le\eta/2\), then \(G\) has no support point in \(\tau+A\).
\end{lemma}

\begin{proof}
Every likelihood ratio is nonnegative, so \(m_G(i)\ge1-W\). Substitution in
the definition of \(\Psi_G\) gives the first inequality, and
\((1-\eta)/(1-\eta/2)<1\) gives the second. \qedhere
\end{proof}

An NPMLE places almost all its mass near the true score origin. The following
coarse estimate isolates a central component before the dimension-specific
edge analysis. For a mixing measure \(G\) in score coordinates, put
\[
    m_G(z)=\int
       \exp\left\{t\cdot z-\frac{\gamma_n|t|^2}{2}\right\}G(dt),
    \qquad
    \mathfrak D_n(G)=\int\{1-e^{-\kappa_n|t|^2}\}\,G(dt).
\]
If \(P\) is standard Gaussian measure, then
\(Pm_G=1-\mathfrak D_n(G)\), and hence
\begin{equation}
\label{eq:support-population-defect}
    P\log m_G\le\log Pm_G
       =\log\{1-\mathfrak D_n(G)\}\le-\mathfrak D_n(G).
\end{equation}

The population loss yields a uniform empirical bound whose coarse
polylogarithmic precision is already sharper than the critical KKT gaps used
in the fixed-dimensional support argument.
\begin{lemma}
\label{lem:support-coarse-defect}
For each fixed \(d\), there is \(C_d<\infty\) such that, uniformly over
\(0\le\eps_n<1\),
\[
    \mathfrak D_n(\widehat G_n)
    =O_{\Pbb}\left\{\frac{(\log n)^{C_d}}{n}\right\}
\]
for every NPMLE \(\widehat G_n\).
\end{lemma}

\begin{proof}
It suffices to optimize over mixing measures supported in
\(\operatorname{conv}\{Z_1/\gamma_n,\ldots,Z_n/\gamma_n\}\). Indeed, if
\(p\) is the Euclidean projection of \(t\) onto this convex hull, then
\[
 (t-p)\cdot Z_i-\frac{\gamma_n}{2}(|t|^2-|p|^2)
 =(t-p)\cdot(Z_i-\gamma_np)-\frac{\gamma_n}{2}|t-p|^2\le0,
\]
so moving \(t\) to \(p\) weakly increases every fitted density. With
probability tending to one this hull lies in
\(B(0,C\sqrt{\log n})\).

For \(\mathcal M(R)=\{m_G:\supp(G)\subset B(0,R)\}\), let
\(N_{[]}(\eta,\mathcal M(R),L^1(P))\) be the least number of brackets
\([l,u]=\{f:l\le f\le u\}\), with \(P(u-l)\le\eta\), needed to cover
\(\mathcal M(R)\). We use the compact Gaussian-mixture entropy bound from
\cite[Section~2]{supportDivergenceCompanion}; see also
\cite[Section~3]{ghosal2001entropies}:
\begin{equation}
\label{eq:support-coarse-entropy}
 \log N_{[]}\{\eta,\mathcal M(R),L^1(P)\}
 \le C_d(1+R)^d\{1+\log(1/\eta)\}^{d+1}.
\end{equation}
The constant is uniform in \(\kappa_n\), since
\[
    m_G(z)\phi_d(z)
    =\int e^{-\kappa_n|t|^2}\phi_d(z-t)\,G(dt),
\]
and an \(L^1(P)\) bracket for \(m_G\) is an ordinary \(L^1\) bracket for
a Gaussian location mixture of total mass at most one.

Write \(\mathbb P_nf=n^{-1}\sum_{i=1}^nf(Z_i)\). Fix a dyadic shell
\(\delta<\mathfrak D_n(G)\le2\delta\), and cover it by upper
brackets \(u\) of \(L^1(P)\)-width \(\delta/4\). If \(m_G\le u\), then
\(Pu\le1-3\delta/4\), and Markov's inequality gives
\[
 \Pbb\{\mathbb P_n\log m_G\ge-\delta/4,\ m_G\le u\}
 \le e^{n\delta/4}(Pu)^n\le e^{-n\delta/2}.
\]
For \(R=O(\{\log n\}^{1/2})\), the logarithm of the number of brackets is
at most \((\log n)^{C_d}\) whenever \(\delta\ge n^{-2}\). A union bound over
brackets and dyadic shells therefore shows that every measure with
\(\mathfrak D_n(G)>C(\log n)^{C_d}/n\) has negative likelihood gain relative to
\(\delta_0\), with probability tending to one. Such a measure cannot be an
NPMLE, proving the claim.
\end{proof}

The defect bound does not rule out several nearby atoms. The local collapse
lemma shows that every maximizer carries all central mass on one atom, with
the remaining edge mass smaller than the KKT gaps used later.
\begin{lemma}
\label{lem:support-fixed-core-collapse}
Assume either \(d=2\) and \(\lambda_n\to\lambda\in(0,\infty)\), or fixed
\(d\ge3\) with \(s_n^\star\) converging in \(\R\). For a
sufficiently small fixed $r_0>0$, every NPMLE can, with
probability tending to one, be written
\begin{equation}
\label{eq:support-central-decomposition}
    \widehat G_n=(1-W_n)\delta_{\tau_n}+G_n^{\rm out},
    \qquad
    \supp(G_n^{\rm out})\subset B(0,r_0)^c,
\end{equation}
where
\[
    W_n+|\tau_n|^2
    =O_{\Pbb}\left\{\frac{(\log n)^{C_d}}{n\kappa_n}\right\}.
\]
In all these regimes, \(W_n=o_{\Pbb}(\kappa_n)\).
\end{lemma}

\begin{proof}
Lemma~\ref{lem:support-coarse-defect} and
\(1-e^{-\kappa_n|t|^2}\ge c\kappa_n\min\{|t|^2,1\}\) imply the displayed
bounds for the mass outside \(B(0,r_0)\) and for the second moment of the
restriction to that ball. Normalize the central restriction to a probability
measure \(G^{\rm cen}\), and let \(\tau\) and \(V\) be its mean and
covariance. Cauchy--Schwarz gives the asserted bound for \(|\tau|^2\).

We show that replacing \(G^{\rm cen}\) by \(\delta_\tau\) strictly increases
the likelihood unless \(V=0\). Keep the outside part fixed and interpolate
\[
 G_u=(1-W)\{(1-u)\delta_\tau+uG^{\rm cen}\}+G^{\rm out},
 \qquad 0\le u\le1.
\]
By concavity it is enough to show that the derivative at \(u=0\) is negative.
With \(h=t-\tau\), put
\[
 g_i(h)=\exp\left\{h\cdot(Z_i-\gamma_n\tau)
                  -\frac{\gamma_n|h|^2}{2}\right\},
 \qquad
 Q_i^{\rm out}=\int q_{n,\tau,t-\tau}(i)\,G^{\rm out}(dt).
\]
The likelihood ratio under \(G_0\) is
\(m_i^0=1-W+Q_i^{\rm out}\). Define
\(\omega_i=(1-W)/m_i^0\). The derivative is
\begin{equation}
\label{eq:support-collapse-derivative}
 \sum_{i=1}^n\omega_i
 \int\{g_i(h)-1\}\,G^{\rm cen}(dt).
\end{equation}

First replace \(\omega_i\) by one. Let
\(\widehat F_\tau(h)=n^{-1}\sum_i g_i(h)\). Uniformly for
\(|\tau|=o_{\Pbb}(\sqrt{\kappa_n})\) and \(|h|\le2r_0\), a compact
finite-dimensional empirical-process bound after two differentiations gives
\[
 \sup_{|h|\le2r_0}
 \|\nabla^2\widehat F_\tau(h)-\nabla^2F_\tau(h)\|_{\rm op}
 =O_{\Pbb}(n^{-1/2})=o_{\Pbb}(\kappa_n).
\]
Here the differentiated class has an integrable envelope of the form
\(C(1+|Z|^2)e^{2r_0|Z|}\), and the population average is
\[
    F_\tau(h)=\exp\{-\kappa_n|h|^2-\gamma_n\tau\cdot h\}.
\]
Indeed,
\[
 \nabla^2F_\tau(h)
 =F_\tau(h)\left[
   (2\kappa_nh+\gamma_n\tau)(2\kappa_nh+\gamma_n\tau)^\top
   -2\kappa_n I_d\right].
\]
The bound on \(|\tau|^2\), together with
\(\kappa_n\gtrsim1/\log n\), gives
\(|\tau|^2=o_{\Pbb}(\kappa_n)\). Thus, after decreasing the fixed \(r_0\) if
needed, \(\nabla^2\widehat F_\tau(h)\preceq-c\kappa_nI_d\) throughout the
ball, with probability tending to one. Taylor's theorem and
\(\int h\,G^{\rm cen}(dt)=0\) now show that the unweighted version of
\eqref{eq:support-collapse-derivative} is at most
\(-cn\kappa_n\operatorname{tr}V\).

It remains to restore the weights. Let \(m_i^1\) be the likelihood ratio
under the original NPMLE \(G_1\), and let
\[
 r_i^1=\frac{\int q_{n,\tau,t-\tau}(i)\,G^{\rm out}(dt)}{m_i^1}
\]
be its posterior responsibility for the outside component. Integrating the
componentwise KKT equality over \(G^{\rm out}\) gives
\(\sum_i r_i^1=nW\).

We next compare the two sets of weights explicitly. Put
\(\bar g_i=\int g_i(h)\,G^{\rm cen}(dt)\), so that
\(m_i^1=(1-W)\bar g_i+Q_i^{\rm out}\). On the event
\(R_n:=\max_i|Z_i|\le C\sqrt{\log n}\),
\[
 e^{-K_n}\le \bar g_i\le e^{K_n},
 \qquad K_n=C_{r_0}(1+R_n),
\]
uniformly in \(i\). Consequently \(m_i^1/m_i^0\le e^{K_n}\), and hence
\[
 1-\omega_i=\frac{Q_i^{\rm out}}{m_i^0}
 \le e^{K_n}r_i^1,
 \qquad
 \sum_i(1-\omega_i)\le e^{K_n}nW.
\]
For each \(i\), another second-order expansion, again using
\(\int h\,G^{\rm cen}(dt)=0\), gives
\[
 \left|\int\{g_i(h)-1\}\,G^{\rm cen}(dt)\right|
 \le C_{r_0}(1+|Z_i|^2)e^{C_{r_0}|Z_i|}\operatorname{tr}V.
\]
It follows that restoring the weights changes
\eqref{eq:support-collapse-derivative} by at most
\[
 C_{r_0}nW(1+\log n)e^{C_{r_0}\sqrt{\log n}}\operatorname{tr}V.
\]
Relative to \(n\kappa_n\operatorname{tr}V\), the last expression is
\[
 O_{\Pbb}\left\{
   \frac{(\log n)^{C_d+1}e^{C_{r_0}\sqrt{\log n}}}
        {n\kappa_n^2}
 \right\}=o_{\Pbb}(1),
\]
since \(\kappa_n\gtrsim1/\log n\) in both stated critical windows. The
derivative is therefore negative whenever \(V\ne0\), contradicting the
maximality of \(G_1\). Hence \(V=0\), which proves
\eqref{eq:support-central-decomposition}. The same rate and
\(\kappa_n\gtrsim1/\log n\) give \(W=o_{\Pbb}(\kappa_n)\).
\end{proof}
